% !TEX root = ../risk_vs.tex

% ==============================
\section{The Most Renowned Research vs Data Mining}\label{sec:app-most-famous}

Section \ref{sec:hetero} showed that journal ranking has relatively little effect. But perhaps one needs to go beyond journal ranking, and focus on the very best research, to find notable outperformance relative to data mining.

To examine this possibility, we take a closer look at the predictability studied in \citet{fama1992cross}, \citet{jegadeesh1993returns}, and \citet{banz1981relationship}. These papers are renowned for studying the predictive power of B/M, momentum, and size, respectively. These findings are not only among the most renowned, but are arguably the ones with the strongest supporting evidence, both theoretical and empirical.  

Theoretical foundations for B/M include \citet{berk1999optimal}, \citet{gomes2003equilibrium}, \citet{campbell2004bad}, \citet{zhang2005value}, and \citet{lettau2007long}. Many of these papers also provide a theoretical foundation for size (see also \citealt{berk1995critique}). Theoretical foundations for momentum include \citet{hong1999unified}, \citet{brav2002competing}, \citet{holden2002news}, and \citet{da2014frog}.\footnote{% 
Other theoretical foundations for B/M and size include \citet{gabaix2008variable}, \citet{papanikolaou2011investment}, and \citet{chen2018general}. For other theoretical foundations for momentum, see \citealt{subrahmanyam2018equity}.
} % end footnote
Many of these theories are themselves award-winning and renowned papers.

Almost all of these theories provide equilibrium foundations---that is, stable relationships between firm characteristics and expected returns. One might argue that the behavioral equilibria are unstable, but others will argue that psychological biases are fundamental, as are limits to arbitrage. And while the multiplicity of theories could be viewed as ``model dredging,'' it could also be viewed as robustness. Theoretical robustness is a feature of physics and statistics, in which core phenomena can be derived from multiple perspectives.\footnote{% begin footnote
Thermodynamic phenomena (e.g. ideal gas law) can be derived from the laws of thermodynamics,  classical mechanics, or quantum mechanics. Core statistical formulas (e.g. regression coefficients), can be derived from method of moments, maximum likelihood, Bayesian assumptions, or data fitting. 
} % end footnote

These renowned papers provide robust empirical evidence. Indeed, \citet{fama1992cross} is in essence a robustness check on \citet{stattman1980book} and \citet{banz1981relationship}. But other empirical papers provide even more robustness (e.g., \citealt{fama1993common}; \citealt{lakonishok1994contrarian}; \citealt{chan1996momentum}; \citealt{asness2013value}). 

Tables \ref{tab:closer_bmdec}-\ref{tab:closer_size} compare these renowned findings to data-mined alternatives. The alternatives come from searching the 29,000 accounting ratios for t-stats and mean returns that are within 10\% and 30\% of the original findings. This filtering ensures the alternatives are similar to the original findings in terms of statistical support. As before,  data-mined t-stats and mean returns use the original papers' stock weighting and sample periods.

\begin{table}[p]
\caption{Data-Mined Predictors that Performed Similarly to Fama-French's B/M (1992)}
\label{tab:closer_bmdec}

\begin{singlespace}
\noindent Table lists 20 of the 163 data-mined predictors that performed similarly to Fama and French's \citeyearpar{fama1992cross} B/M in the original sample period. It includes predictors with t-stats within 10\% and mean returns within 30\% of the original findings. Signals are ranked by the absolute difference in mean return. Sign = -1 indicates that a high signal implies a lower mean return in-sample. Data mining performs similarly to trading on \citepos{fama1992cross} B/M.
\end{singlespace}
\begin{centering}
\vspace{0ex}
\par\end{centering}
\centering{}\setlength{\tabcolsep}{0.7ex} \small
\begin{center}
% ==== begin paste
% Table generated by Excel2LaTeX from sheet 'inspect BMdec' 
\begin{tabular}{rlrrr} \toprule 
\multicolumn{1}{c}{Similarity} 
    & \multirow{2}[2]{*}{Signal} 
    & \multicolumn{1}{l}{\multirow{2}[2]{*}{Sign}} & \multicolumn{2}{c}{Mean Return (\% p.m.)} \\   
\multicolumn{1}{c}{Rank} 
    &   &   
    & \multicolumn{1}{c}{1963-1990} & \multicolumn{1}{c}{1991-2023} \\ \midrule \multicolumn{5}{l}{\textit{Peer-Reviewed}} \\ 
\midrule   
% latex table generated in R 4.2.3 by xtable 1.8-4 package
% Mon Mar  3 17:42:07 2025
  & Book / Market (Fama-French 1992) & 1 & 0.96 & 0.61 \\ 
  \midrule 
\multicolumn{5}{l}{\textit{Data-Mined}} \\ 
\midrule 
 1 & $\Delta$[PPE net]/lag[Sales] & -1 & 0.96 & 0.73 \\ 
  2 & $\Delta$[Assets]/lag[Cost of goods sold] & -1 & 0.95 & 0.80 \\ 
  3 & $\Delta$[Assets]/lag[Operating expenses] & -1 & 0.95 & 0.84 \\ 
  4 & [Depreciation (CF acct)]/[Capex PPE sch V] & 1 & 0.97 & 0.68 \\ 
  5 & [Stock issuance]/[Debt in current liab] & -1 & 0.94 & 0.73 \\ 
  6 & $\Delta$[Assets]/lag[SG\&A] & -1 & 0.94 & 0.78 \\ 
  7 & $\Delta$[PPE net]/lag[Gross profit] & -1 & 0.98 & 0.45 \\ 
  8 & $\Delta$[PPE net]/lag[Current liabilities] & -1 & 0.94 & 0.85 \\ 
  9 & [Stock issuance]/[Capex PPE sch V] & -1 & 0.94 & 1.00 \\ 
  10 & $\Delta$[PPE (gross)]/lag[Gross profit] & -1 & 0.93 & 0.33 \\ 
   & . . . &  &  &  \\ 
  101 & $\Delta$[Assets]/lag[Assets other sundry] & -1 & 0.75 & 0.95 \\ 
  102 & $\Delta$[Liabilities]/lag[Invest tax credit inc ac] & -1 & 0.74 & 0.14 \\ 
  103 & $\Delta$[PPE net]/lag[Capital expenditure] & -1 & 0.74 & 0.79 \\ 
  104 & $\Delta$[PPE net]/lag[Interest expense] & -1 & 0.75 & 0.63 \\ 
  105 & $\Delta$[Receivables]/lag[Assets] & -1 & 0.74 & 0.59 \\ 
   & . . . &  &  &  \\ 
  159 & $\Delta$[Assets]/lag[IB adjusted for common s] & -1 & 0.67 & -0.02 \\ 
  160 & $\Delta$[Assets]/lag[Income bf extraordinary] & -1 & 0.67 & -0.03 \\ 
  161 & $\Delta$[Assets]/lag[Net income] & -1 & 0.67 & -0.01 \\ 
  162 & $\Delta$[Cost of goods sold]/lag[Current liabilities] & -1 & 0.67 & 0.65 \\ 
  163 & $\Delta$[Inventories]/lag[Curr assets other sundry] & -1 & 0.67 & 0.63 \\ 
  \midrule 
  & Mean Data-Mined &  & 0.83 & 0.65 \\ 
  
\bottomrule \end{tabular}% 
% ==== end paste
\end{center} 
\end{table}

\begin{table}[p]
\caption{Data-Mined Predictors That Performed Similarly to Jegadeesh and Titman's
12-Month Momentum (1993) }
\label{tab:closer_mom12m}

\begin{singlespace}
\noindent Table lists 20 of the 44 data-mined predictors that performed similarly to Jegadeesh and Titman's \citeyearpar{jegadeesh1993returns} 12-month momentum in the original sample period. It includes predictors with t-stats within 10\% and mean returns within 30\% of the original findings. Signals are ranked by the absolute difference in mean return. Sign = -1 indicates that a high signal implies a lower mean return in-sample. Data mining somewhat underperforms \citepos{jegadeesh1993returns} momentum.
\end{singlespace}
\begin{centering}
\vspace{0ex}
\par\end{centering}
\centering{}\setlength{\tabcolsep}{0.5ex} \small
\begin{center}
% ==== begin paste
% Table generated by Excel2LaTeX from sheet 'inspect Mom12m' 
\begin{tabular}{rlrrr} \toprule 
\multicolumn{1}{c}{Similarity} 
    & \multirow{2}[2]{*}{Signal} 
    & \multicolumn{1}{l}{\multirow{2}[2]{*}{Sign}} & \multicolumn{2}{c}{Mean Return (\% p.m.)} \\   
\multicolumn{1}{c}{Rank}  &  &   & \multicolumn{1}{c}{1964-1989} & \multicolumn{1}{c}{1990-2023} \\ 
\midrule 
\multicolumn{5}{l}{\textit{Peer-Reviewed}} \\ 
\midrule   
% === insert 
% latex table generated in R 4.5.3 by xtable 1.8-8 package
% Thu Sep 24 16:05:29 2026
  & 12-Month Momentum (Jegadeesh-Titman 1993) & 1 & 1.36 & 0.72 \\ 
  \midrule 
\multicolumn{5}{l}{\textit{Data-Mined}} \\ 
\midrule 
 1 & [Retained earnings unadj]/[Liabilities other] & 1 & 1.37 & 0.21 \\ 
  2 & [Retained earnings unadj]/[Market equity FYE] & 1 & 1.38 & -0.02 \\ 
  3 & [Retained earnings unadj]/[Assets other sundry] & 1 & 1.40 & 0.20 \\ 
  4 & [Retained earnings unadj]/[Cash \& ST investments] & 1 & 1.42 & 0.31 \\ 
  \midrule 
  & Mean Data-Mined &  & 1.39 & 0.17 \\ 
  
\bottomrule 
\end{tabular}% 
% ==== end paste
\end{center} 
\end{table}


\begin{table}[p]
  \caption{Data-Mined Predictors That Performed Similarly to Banz's Size (1981)}
    \label{tab:closer_size}
    
    \begin{singlespace}
    \noindent Table lists 20 of the 220 data-mined predictors that performed similarly to Banz's \citeyearpar{banz1981relationship} size in the original sample period. It includes predictors with t-stats within 10\% and mean returns within 30\% of the original findings. Signals are ranked by the absolute difference in mean return. Sign = -1 indicates that a high signal implies a lower mean return in-sample. Data mining outperforms \citepos{banz1981relationship} size.
    \end{singlespace}
    \begin{centering}
    \vspace{-1ex}
    \par\end{centering}
    \centering{}\setlength{\tabcolsep}{1.0ex} \footnotesize
    \begin{center}
    % ==== begin paste
    % Table generated by Excel2LaTeX from sheet 'inspect Size' 
    \begin{tabular}{rlrrr} \toprule 
    \multicolumn{1}{c}{Similarity} 
        & \multirow{2}[2]{*}{Signal} 
        & \multicolumn{1}{l}{\multirow{2}[2]{*}{Sign}} & \multicolumn{2}{c}{Mean Return (\% Monthly)} \\   
    \multicolumn{1}{c}{Rank} 
        &  &  & \multicolumn{1}{c}{1926-1975} & \multicolumn{1}{c}{1976-2023} \\ \midrule \multicolumn{5}{l}{\textit{Peer-Reviewed}} \\ \midrule   
    % latex table generated in R 4.2.3 by xtable 1.8-4 package
% Mon Mar  3 17:42:08 2025
  & Size (Banz 1981) & -1 & 0.50 & 0.15 \\ 
  \midrule 
\multicolumn{5}{l}{\textit{Data-Mined}} \\ 
\midrule 
 1 & $\Delta$[Equity liq value]/lag[Sales] & -1 & 0.50 & 0.72 \\ 
  2 & [Invested capital]/[Market equity FYE] & 1 & 0.50 & 0.83 \\ 
  3 & $\Delta$[Assets]/lag[Pref stock liq value] & -1 & 0.49 & 0.18 \\ 
  4 & $\Delta$[Equity liq value]/lag[Current liabilities] & -1 & 0.48 & 0.79 \\ 
  5 & $\Delta$[Receivables]/lag[Pref stock redemp val] & -1 & 0.48 & 0.10 \\ 
  6 & $\Delta$[Current assets]/lag[Invest tax credit inc ac] & -1 & 0.52 & 0.35 \\ 
  7 & $\Delta$[Assets]/lag[Pref stock redemp val] & -1 & 0.47 & 0.23 \\ 
  8 & $\Delta$[Equity liq value]/lag[Curr assets other sundry] & -1 & 0.48 & 0.69 \\ 
  9 & $\Delta$[Common equity tangible]/lag[SG\&A] & -1 & 0.47 & 0.40 \\ 
  10 & $\Delta$[Invested capital]/lag[PPE (gross)] & -1 & 0.47 & 0.90 \\ 
   & . . . &  &  &  \\ 
  101 & $\Delta$[Depreciation \& amort]/lag[Common equity tangible] & -1 & 0.39 & 0.40 \\ 
  102 & $\Delta$[Depreciation \& amort]/lag[Invest \& advances other] & -1 & 0.38 & 0.52 \\ 
  103 & $\Delta$[Depreciation depl amort]/lag[Interest expense] & -1 & 0.39 & 0.07 \\ 
  104 & $\Delta$[Num employees]/lag[Long-term debt] & -1 & 0.39 & 0.55 \\ 
  105 & $\Delta$[Num employees]/lag[Invest \& advances other] & -1 & 0.39 & 0.45 \\ 
   & . . . &  &  &  \\ 
  216 & $\Delta$[Pref stock nonredeemable]/lag[PPE (gross)] & -1 & 0.35 & 0.69 \\ 
  217 & $\Delta$[Receivables]/lag[Curr assets other sundry] & -1 & 0.35 & 0.60 \\ 
  218 & $\Delta$[Operating expenses]/lag[Invested capital] & -1 & 0.35 & 0.62 \\ 
  219 & [Acquisitions]/[Nonop income] & -1 & 0.65 & 0.15 \\ 
  220 & [Acquisitions]/[Operating expenses] & -1 & 0.64 & 0.34 \\ 
  \midrule 
  & Mean Data-Mined &  & 0.44 & 0.42 \\ 
  
    \bottomrule \end{tabular}% 
    % ==== end paste
    \end{center} 
\end{table}


% walk through one table
Table \ref{tab:closer_bmdec} applies this exercise to \citepos{fama1992cross} B/M. It lists 20 of the 163 data-mined predictors that performed similarly to B/M in \citepos{fama1992cross} 1963-1990 sample period. The predictors are sorted by the absolute difference in the original sample mean return. By this metric, the most similar predictor to B/M is $\Delta \text{[PPE net]} / \text{lag [Sales]}$, which can be thought of as a measure of investment. This predictor earned 96 bps per month in the original sample period, identical to B/M. The `1991-2023' column shows that $\Delta \text{[PPE net]} / \text{lag [Sales]}$ slightly outperforms B/M post-sample, earning 73 bps per month compared to 61 bps for B/M. 

Other data-mined alternatives to B/M include those related to equity issuance ([Stock issuance]/ [Debt in current liab]) and accruals ($\Delta$ [Receivables]/ lag [Assets]). The post-sample performance of these data-mined alternatives varies. But on average, they are quite similar to \citet{fama1992cross}. Trading on Fama and French's finding would have earned 61 bps per month post-sample, compared to 65 bps for the typical data mined alternative.

Table \ref{tab:closer_mom12m} applies the same exercise to \citepos{jegadeesh1993returns} 12-month momentum. Since momentum has a much higher mean return, only 44 data-mined alternatives are found. Here, data mining underperforms on average, earning 52 bps compared to 72 bps for \citepos{jegadeesh1993returns} finding. However, Table \ref{tab:closer_size} shows data mining outperforming, earning 42 bps compared to 15 bps for \citepos{banz1981relationship} size.

Averaging across the three tables, data mining performs similarly to these renowned findings. Though the samples are small, they suggest that focusing on the best research does not significantly affect our results.  